% Chapter: Cubical Structure and the Path Algebra
% Part of the CCTT Monograph — Proof Calculus, Chapter II

\chapter{Cubical Structure and the Path Algebra}
\label{ch:c4-cubical}

In Chapter~\ref{ch:c4-calculus} we introduced the syntax and typing rules of
$\mathrm{C}^4$.  This chapter develops the \emph{cubical} fragment in depth:
the interval, connections, Kan operations, and the derived principles of
function extensionality and univalence.  We then prove the central theorem
linking cubical path composition to CCTT axiom chain composition, establishing
that the proof-theoretic machinery of Part~IX is faithfully embedded in the
type-theoretic framework.


% ─────────────────────────────────────────────────────────
\section{The Interval and Its Connection Algebra}
\label{sec:interval}

\begin{definition}[The Interval]
The \textbf{interval} $\interval$ is a type with two distinguished points
$\mathsf{0}_\interval, \mathsf{1}_\interval : \interval$ and three
connection operations:
\begin{align}
  \_ \wedge \_ &: \interval \times \interval \to \interval
    \tag{meet / min} \\
  \_ \vee \_ &: \interval \times \interval \to \interval
    \tag{join / max} \\
  \sim &: \interval \to \interval
    \tag{reversal}
\end{align}
subject to the following equalities, which hold definitionally:
\begin{enumerate}
\item $\sim \mathsf{0} = \mathsf{1}$, \quad $\sim \mathsf{1} = \mathsf{0}$,
  \quad $\sim(\sim r) = r$.
\item $r \wedge \mathsf{0} = \mathsf{0}$, \quad $r \wedge \mathsf{1} = r$,
  \quad $r \vee \mathsf{0} = r$, \quad $r \vee \mathsf{1} = \mathsf{1}$.
\item $r \wedge s = s \wedge r$, \quad $r \vee s = s \vee r$.
\item $(r \wedge s) \wedge t = r \wedge (s \wedge t)$, \quad
  $(r \vee s) \vee t = r \vee (s \vee t)$.
\item De Morgan: $\sim(r \wedge s) = (\sim r) \vee (\sim s)$, \quad
  $\sim(r \vee s) = (\sim r) \wedge (\sim s)$.
\end{enumerate}
\end{definition}

\begin{remark}
The connection algebra $(\interval, \wedge, \vee, \sim, \mathsf{0},
\mathsf{1})$ forms a De Morgan algebra.  This is the same structure used by
Cohen--Coquand--Huber--M\"ortberg (CCHM)~\cite{cohen2018}.  The key property
is that $\interval$ is \emph{not} a type with decidable equality---there is
no function $\interval \to \mathsf{2}$ separating $\mathsf{0}$ from
$\mathsf{1}$.
\end{remark}

\begin{definition}[Face Lattice]
A \textbf{face formula} $\varphi$ is a finite conjunction/disjunction of
atomic constraints:
\[
  \varphi, \psi ::= (r = \mathsf{0}) \mid (r = \mathsf{1})
    \mid \varphi \wedge \psi \mid \varphi \vee \psi \mid \mathsf{1}_\mathbb{F}
    \mid \mathsf{0}_\mathbb{F}
\]
We write $\mathbb{F}$ for the \textbf{face lattice} of all face formulas.
A context $\Gamma, \varphi$ extends $\Gamma$ with the constraint $\varphi$,
restricting to the face where $\varphi$ holds.
\end{definition}

\begin{definition}[Partial Elements]
Given $\Gamma \vdash A : \mathcal{U}_i$ and a face formula $\varphi$, a
\textbf{partial element} of $A$ along $\varphi$ is a term
$\Gamma, \varphi \vdash u : A$.  We write this as
$u : A\,[\varphi]$ and call $\varphi$ the \textbf{extent} of $u$.
\end{definition}


% ─────────────────────────────────────────────────────────
\section{Path Types in Detail}
\label{sec:path-types}

\begin{definition}[Path Type]
For $A : \mathcal{U}_i$ and $a, b : A$, the \textbf{path type} is:
\[
  \Path_A\, a\, b \;\defeq\;
  \{p : \interval \to A \mid p\,\mathsf{0} \equiv a,\;
    p\,\mathsf{1} \equiv b\}
\]
\end{definition}

\noindent
Path types internalize propositional equality as continuous deformation.
The typing rules from Chapter~\ref{ch:c4-calculus} give:

\begin{gather}
\frac{\Gamma, i : \interval \vdash t : A \qquad
      t[i := \mathsf{0}] \equiv a \qquad t[i := \mathsf{1}] \equiv b}
     {\Gamma \vdash \langle i \rangle\, t : \Path_A\, a\, b}
\tag{PathI}
\\[6pt]
\frac{\Gamma \vdash p : \Path_A\, a\, b \qquad
      \Gamma \vdash r : \interval}
     {\Gamma \vdash p\; r : A}
\tag{PathE}
\end{gather}

with the definitional equalities:
\begin{align}
  (\langle i \rangle\, t)\; r &\equiv t[i := r]
    & \text{($\beta$ for paths)} \tag{Path-$\beta$} \\
  \langle i \rangle\, (p\; i) &\equiv p
    & \text{($\eta$ for paths, when $i \notin \mathrm{FV}(p)$)}
    \tag{Path-$\eta$}
\end{align}

\begin{definition}[Path Operations]
\label{def:path-ops}
Using the connection algebra, we define:
\begin{enumerate}
\item \textbf{Constant path (reflexivity):}
  $\mathsf{refl}_a \defeq \langle i \rangle\, a : \Path_A\, a\, a$.

\item \textbf{Symmetry (path reversal):}
  Given $p : \Path_A\, a\, b$, define
  $p^{-1} \defeq \langle i \rangle\, p\,(\sim i) : \Path_A\, b\, a$.

\item \textbf{Connection terms:}
  Given $p : \Path_A\, a\, b$, the \textbf{connections} are:
  \begin{align*}
    p^{\wedge} &\defeq \langle i \rangle\langle j \rangle\,
      p\,(i \wedge j) &&: \Path_{\Path_A\, a\, b}\,
      \mathsf{refl}_a\, p \\
    p^{\vee} &\defeq \langle i \rangle\langle j \rangle\,
      p\,(i \vee j) &&: \Path_{\Path_A\, a\, b}\,
      p\, \mathsf{refl}_b
  \end{align*}
\end{enumerate}
\end{definition}

\begin{proposition}[Groupoid Laws]
\label{prop:groupoid}
For paths $p : \Path_A\, a\, b$, $q : \Path_A\, b\, c$,
$r : \Path_A\, c\, d$, the following hold (as paths of paths):
\begin{enumerate}
\item $p \cdot \mathsf{refl}_b \sim p$ \quad (right unit)
\item $\mathsf{refl}_a \cdot p \sim p$ \quad (left unit)
\item $p \cdot p^{-1} \sim \mathsf{refl}_a$ \quad (right inverse)
\item $(p \cdot q) \cdot r \sim p \cdot (q \cdot r)$ \quad (associativity)
\end{enumerate}
where $p \cdot q$ denotes path composition defined via $\hfilling$.
\end{proposition}

\begin{proof}
Each identity is constructed using the connection algebra.  For (1),
the term $\langle i\rangle\langle j\rangle\, p\,(i \vee j)$ provides
a path from $p \cdot \mathsf{refl}_b$ to $p$: at $j = \mathsf{0}$ we
get $p \cdot \mathsf{refl}_b$ (by the definition of composition using
$\hfilling$) and at $j = \mathsf{1}$ we get $\langle i \rangle\, p\,i = p$.
The remaining identities follow by analogous two-dimensional constructions
using $\wedge$, $\vee$, and $\sim$.
\end{proof}


% ─────────────────────────────────────────────────────────
\section{Kan Operations}
\label{sec:kan-ops}

The Kan operations give path types their computational content.

\begin{definition}[Homogeneous Composition --- $\hfilling$]
\label{def:hcomp}
Given $A : \mathcal{U}_i$, a face formula $\varphi$, a partial path
$\Gamma, \varphi, i : \interval \vdash u : A$, and a base point
$\Gamma \vdash u_0 : A$ agreeing with $u$ at $i = \mathsf{0}$
(i.e., $u[\mathsf{0}/i] \equiv u_0$ when $\varphi$ holds):
\[
\frac{\Gamma \vdash A : \mathcal{U}_i \qquad
      \Gamma, \varphi, i : \interval \vdash u : A \qquad
      \Gamma \vdash u_0 : A\,[\varphi \mapsto u[\mathsf{0}/i]]}
     {\Gamma \vdash \hfilling^i\, A\, [\varphi \mapsto u]\, u_0 : A}
\tag{hcomp}
\]
with the boundary law: when $\varphi$ holds,
$\hfilling^i\, A\, [\varphi \mapsto u]\, u_0 \equiv u[\mathsf{1}/i]$.
\end{definition}

\begin{definition}[Transport --- $\transp$]
\label{def:c4-transp}
Given a line of types $\Gamma, i : \interval \vdash A : \mathcal{U}_j$
and a term $\Gamma \vdash t : A[\mathsf{0}/i]$:
\[
\frac{\Gamma, i : \interval \vdash A : \mathcal{U}_j \qquad
      \Gamma \vdash t : A[\mathsf{0}/i]}
     {\Gamma \vdash \transp^i\, A\, t : A[\mathsf{1}/i]}
\tag{transp}
\]
When $A$ is constant in $i$ (i.e., $i \notin \mathrm{FV}(A)$), transport
is the identity: $\transp^i\, A\, t \equiv t$.
\end{definition}

\begin{definition}[Path Composition via $\hfilling$]
\label{def:path-comp}
Given $p : \Path_A\, a\, b$ and $q : \Path_A\, b\, c$, their
\textbf{composition} $p \cdot q : \Path_A\, a\, c$ is:
\[
  p \cdot q \;\defeq\; \langle i \rangle\,
    \hfilling^j\, A\,
    [(i = \mathsf{0}) \mapsto a,\;
     (i = \mathsf{1}) \mapsto q\,j]\,
    (p\,i)
\]
where $j$ is fresh.  At $i = \mathsf{0}$ this gives $a$; at
$i = \mathsf{1}$ it gives $\hfilling^j\, A\, [\cdots]\, b$ which
computes to $q\,\mathsf{1} = c$.
\end{definition}

\noindent
\textbf{Computation rules for Kan operations on type formers.}
The Kan operations compute structurally on each type former:

\begin{enumerate}
\item \textbf{$\Pi$-types:}
  $\transp^i\, (\Pi(x:A).\,B)\, f \equiv
   \lambda(x : A[\mathsf{1}/i]).\,
   \transp^i\, B[x := \transp^i\, (\sim A)\, x]\,
   (f\, (\transp^i\, (\sim A)\, x))$

\item \textbf{$\Sigma$-types:}
  $\transp^i\, (\Sigma(x:A).\,B)\, (a, b) \equiv
   (a', \transp^i\, B[x := p]\, b)$
  where $a' = \transp^i\, A\, a$ and $p$ is the transport path.

\item \textbf{Path types:}
  Transport in $\Path_A\, a\, b$ composes paths using $\hfilling$.

\item \textbf{Natural numbers:} $\transp^i\, \Nat\, n \equiv n$
  (discrete type).

\item \textbf{Inductive types:} Transport applies constructor-wise,
  preserving the recursive structure.
\end{enumerate}


% ─────────────────────────────────────────────────────────
\section{Function Extensionality}
\label{sec:c4-funext}

\begin{theorem}[Function Extensionality]
\label{thm:c4-funext}
If $f, g : \Pi(x : A).\, B(x)$ and
$h : \Pi(x : A).\, \Path_{B(x)}\, (f\,x)\, (g\,x)$, then
there exists a path $\mathsf{funext}(h) : \Path_{\Pi(x:A).\,B(x)}\, f\, g$.
\end{theorem}

\begin{proof}
Define $\mathsf{funext}(h) \defeq \langle i \rangle\, \lambda(x : A).\,
h\,x\,i$.  We check the boundary:
\begin{itemize}
\item At $i = \mathsf{0}$:
  $(\langle i \rangle\, \lambda x.\, h\,x\,i)\,\mathsf{0}
   = \lambda x.\, h\,x\,\mathsf{0}
   = \lambda x.\, f\,x = f$.
\item At $i = \mathsf{1}$:
  $(\langle i \rangle\, \lambda x.\, h\,x\,i)\,\mathsf{1}
   = \lambda x.\, h\,x\,\mathsf{1}
   = \lambda x.\, g\,x = g$.
\end{itemize}
Hence $\mathsf{funext}(h) : \Path\, f\, g$.  This is a
\emph{definitional} proof---no axiom is needed.
\end{proof}

\begin{remark}
Function extensionality is a theorem in CIC only with the axiom of
functional extensionality (Lean's \texttt{funext}).  In $\mathrm{C}^4$,
it is \emph{derived} from the cubical structure.  This means
$\mathrm{C}^4$ is strictly more powerful than CIC without funext,
while remaining conservative over CIC + funext.
\end{remark}


% ─────────────────────────────────────────────────────────
\section{Glue Types and Univalence}
\label{sec:univalence}

\begin{definition}[Equivalence]
A function $f : A \to B$ is an \textbf{equivalence} if it has
contractible fibers:
\[
  \mathsf{isEquiv}(f) \;\defeq\;
  \Pi(b : B).\, \mathsf{isContr}\bigl(
    \Sigma(a : A).\, \Path_B\, (f\,a)\, b
  \bigr)
\]
We write $A \simeq B$ for the type of equivalences $\Sigma(f : A \to B).\,
\mathsf{isEquiv}(f)$.
\end{definition}

\begin{definition}[Glue Types]
\label{def:glue}
Given $\Gamma \vdash A : \mathcal{U}_i$, a face formula $\varphi$, and
along $\varphi$ a type $T : \mathcal{U}_i$ with an equivalence
$e : T \simeq A$:
\[
  \frac{\Gamma \vdash A : \mathcal{U}_i \qquad
        \Gamma, \varphi \vdash T : \mathcal{U}_i \qquad
        \Gamma, \varphi \vdash e : T \simeq A}
       {\Gamma \vdash \mathsf{Glue}\,[\varphi \mapsto (T, e)]\, A
         : \mathcal{U}_i}
  \tag{Glue}
\]
with the boundary law: when $\varphi$ holds,
$\mathsf{Glue}\,[\varphi \mapsto (T, e)]\, A \equiv T$.
\end{definition}

\begin{theorem}[Univalence]
\label{thm:univalence}
For any $A, B : \mathcal{U}_i$, the canonical map
\[
  \mathsf{idToEquiv} : \Path_{\mathcal{U}_i}\, A\, B \to (A \simeq B)
\]
defined by $\mathsf{idToEquiv}(p) \defeq \transp^i\, (p\,i)$ together
with its equivalence proof, is itself an equivalence.
\end{theorem}

\begin{proof}[Proof sketch]
We construct the inverse $\mathsf{ua} : (A \simeq B) \to
\Path_{\mathcal{U}_i}\, A\, B$ using Glue types.  Given an equivalence
$e : A \simeq B$, define:
\[
  \mathsf{ua}(e) \defeq \langle i \rangle\,
    \mathsf{Glue}\,[(i = \mathsf{0}) \mapsto (A, e),\;
                     (i = \mathsf{1}) \mapsto (B, \mathsf{id}_B)]\, B
\]
At $i = \mathsf{0}$ this is $A$ (glued to $B$ via $e$); at $i = \mathsf{1}$
this is $B$ (glued to $B$ via identity).  The round-trip
$\mathsf{idToEquiv} \circ \mathsf{ua} \sim \mathrm{id}$ and
$\mathsf{ua} \circ \mathsf{idToEquiv} \sim \mathrm{id}$ are verified
using the computation rules for $\transp$ and $\mathsf{Glue}$.
\end{proof}


% ─────────────────────────────────────────────────────────
\section{Path Composition and Axiom Chain Correspondence}
\label{sec:path-axiom}

This section proves the central theorem connecting the proof-theoretic
machinery of Chapters~\ref{ch:proof-theory-foundations}--\ref{ch:proof-theory-automation}
to the type-theoretic framework of $\mathrm{C}^4$.

\begin{definition}[Axiom Path]
\label{def:axiom-path}
For each CCTT axiom $D_k : \ell_k \equiv r_k$ (with $k \in \{1, \ldots, 24\}$
or from the extended axiom set), the path constructor of the $\PyObj$ HIT
gives a term:
\[
  \mathsf{D}_k : \Pi(\sigma : \mathsf{Subst}).\,
    \Path_{\mathcal{U}_\PyObj}\, \den{\ell_k\sigma}\, \den{r_k\sigma}
\]
We call $\mathsf{D}_k(\sigma)$ an \textbf{axiom path}.
\end{definition}

\begin{definition}[Axiom Chain]
\label{def:axiom-chain}
An \textbf{axiom chain} from $\den{f}$ to $\den{g}$ is a sequence
\[
  \den{f} = t_0 \xrightarrow{D_{i_1}(\sigma_1)} t_1
    \xrightarrow{D_{i_2}(\sigma_2)} t_2
    \xrightarrow{\quad\cdots\quad} t_{n-1}
    \xrightarrow{D_{i_n}(\sigma_n)} t_n = \den{g}
\]
where each step applies a (possibly reversed) axiom at a subterm position.
\end{definition}

\begin{theorem}[Path--Axiom Correspondence]
\label{thm:path-axiom}
Let $f, g$ be Python programs with OTerm denotations $\den{f}, \den{g}
\in \mathcal{U}_\PyObj$.
\begin{enumerate}
\item \textbf{Axiom chain $\Rightarrow$ path:}
  If there is an axiom chain $D_{i_1}(\sigma_1) \cdot D_{i_2}(\sigma_2)
  \cdots D_{i_n}(\sigma_n)$ from $\den{f}$ to $\den{g}$, then there
  exists a term $p : \Path_{\mathcal{U}_\PyObj}\, \den{f}\, \den{g}$
  constructible from the HIT path constructors and cubical path
  composition.

\item \textbf{Path $\Rightarrow$ axiom chain:}
  Conversely, if $p : \Path_{\mathcal{U}_\PyObj}\, \den{f}\, \den{g}$ is
  constructible from the HIT path constructors (without transport through
  non-trivial type families), then there exists a finite axiom chain from
  $\den{f}$ to $\den{g}$.
\end{enumerate}
\end{theorem}

\begin{proof}
\textbf{Part (1).}  We proceed by induction on the length $n$ of the
axiom chain.

\emph{Base case} ($n = 0$): $\den{f} \equiv \den{g}$, so
$\mathsf{refl}_{\den{f}} : \Path\, \den{f}\, \den{g}$.

\emph{Inductive step}: Suppose the chain has $n$ steps, with the
first step $D_{i_1}(\sigma_1) : \den{f} \equiv t_1$ applied at subterm
position $\pi$.  By the HIT structure, $D_{i_1}(\sigma_1)$ gives a
path $p_1 : \Path\, \den{\ell_{i_1}\sigma_1}\, \den{r_{i_1}\sigma_1}$.
To apply this at position $\pi$, we use congruence: if $\den{f} =
C[\ell_{i_1}\sigma_1]$ for a context $C[\cdot]$, then
$\mathsf{ap}_{C}(p_1) : \Path\, C[\ell_{i_1}\sigma_1]\,
C[r_{i_1}\sigma_1] = \Path\, \den{f}\, t_1$.  By the inductive
hypothesis, the remaining chain gives $q : \Path\, t_1\, \den{g}$.
Compose: $\mathsf{ap}_C(p_1) \cdot q : \Path\, \den{f}\, \den{g}$.

\textbf{Part (2).}  A path $p$ in the HIT $\PyObj$ that is built only
from HIT constructors is, by the elimination principle of HITs, a
composition of the generating paths $D_k(\sigma)$ and their inverses
(possibly applied under contexts via $\mathsf{ap}$).  Each such
composition step corresponds to an axiom application.  The elimination
principle gives a term in the free $\infty$-groupoid generated by the
path constructors; truncating to the 0-level (equality of OTerms)
gives a finite word in the axiom applications, i.e., an axiom chain.
\end{proof}

\begin{corollary}
The CCTT proof checker (which searches for axiom chains) is
\textbf{sound and complete} with respect to $\mathrm{C}^4$ path
equality in $\mathcal{U}_\PyObj$, restricted to HIT-constructible
paths.
\end{corollary}

\begin{proof}
Immediate from Theorem~\ref{thm:path-axiom}: the checker finds an
axiom chain if and only if a HIT-constructible path exists.
\end{proof}


% ─────────────────────────────────────────────────────────
\section{Connection to the PyComp HIT}
\label{sec:pycomp-hit}

The $\PyObj$ HIT from Chapter~\ref{ch:c4-calculus} has two classes of
constructors:

\begin{enumerate}
\item \textbf{Point constructors}: $\mathsf{OVar}$, $\mathsf{OLit}$,
  $\mathsf{OOp}$, $\mathsf{OFold}$, $\mathsf{OCase}$, $\mathsf{OFix}$,
  $\mathsf{OApp}$, $\mathsf{OLam}$, $\mathsf{OSeq}$, $\ldots$
  These are the syntax tree nodes of the OTerm language.

\item \textbf{Path constructors}: $\mathsf{D}_1, \ldots, \mathsf{D}_{24}$
  and the extended axioms $\mathsf{P}_1, \ldots, \mathsf{P}_{48}$.
  These identify OTerms that are semantically equal.
\end{enumerate}

\begin{definition}[Path Composition in $\PyObj$]
Given axiom paths $p_k : \Path_\PyObj\, s_{k-1}\, s_k$ for
$k = 1, \ldots, n$, the \textbf{chain composition} is:
\[
  p_1 \cdot p_2 \cdots p_n \;\defeq\;
  \underbrace{(\cdots((p_1 \cdot p_2) \cdot p_3) \cdots) \cdot p_n}_{
    \text{left-associated via Definition~\ref{def:path-comp}}}
  : \Path_\PyObj\, s_0\, s_n
\]
\end{definition}

\begin{example}[Factorial Equivalence as Path]
The proof that recursive factorial equals iterative factorial
(Example~1 in Chapter~\ref{ch:proof-theory-examples}) corresponds to a
path in $\PyObj$:
\begin{align*}
  &\den{\mathsf{fact\_rec}(n)} \\
  &\quad \xrightarrow{\mathsf{NatInd}(\text{base, step})}
    \mathsf{OFold}(\times, 1, \mathsf{range}(1, n{+}1)) \\
  &\quad \xrightarrow{D_{15}(\sigma)}
    \den{\mathsf{fact\_iter}(n)}
\end{align*}
In $\mathrm{C}^4$, this is a path:
\[
  p_{\mathsf{fact}} \defeq
  \mathsf{ap}_{\mathsf{OFix}}(\mathsf{nat\_ind\_path}) \cdot
  \mathsf{D}_{15}(\sigma)
  : \Path_\PyObj\, \den{\mathsf{fact\_rec}}\, \den{\mathsf{fact\_iter}}
\]
\end{example}

\begin{example}[Sum-Range Identity as Path]
The identity $\mathsf{sum}(\mathsf{range}(n)) = n(n-1)/2$ is a
two-step path:
\[
  \den{\mathsf{sum}(\mathsf{range}(n))}
  \xrightarrow{P_{12}}
  \mathsf{OFold}(+, 0, \mathsf{range}(n))
  \xrightarrow{D_{18}}
  \den{n(n-1)/2}
\]
In $\mathrm{C}^4$: $P_{12}(\sigma) \cdot D_{18}(\sigma') :
\Path_\PyObj\, \den{\mathsf{sum}(\mathsf{range}(n))}\,
\den{n(n-1)/2}$.
\end{example}


% ─────────────────────────────────────────────────────────
\section{Cubical Structure Interacting with Fibers}
\label{sec:cubical-fibers}

When the site $\site$ is non-trivial, cubical and cohomological
structure interact:

\begin{proposition}[Restriction Commutes with Path Operations]
\label{prop:restr-path}
For $p : \Path_A\, a\, b$ and $U \in \mathrm{Ob}(\site)$:
\begin{enumerate}
\item $\restr{p}{U} : \Path_{\restr{A}{U}}\, \restr{a}{U}\, \restr{b}{U}$
\item $\restr{(p \cdot q)}{U} \equiv \restr{p}{U} \cdot \restr{q}{U}$
\item $\restr{(p^{-1})}{U} \equiv (\restr{p}{U})^{-1}$
\item $\restr{(\transp^i\, A\, t)}{U} \equiv
  \transp^i\, \restr{A}{U}\, \restr{t}{U}$
\end{enumerate}
\end{proposition}

\begin{proof}
(1) follows from the restriction computation rule (Path-res) in
Chapter~\ref{ch:c4-calculus}: $\restr{(\langle i \rangle\, t)}{U}
\equiv \langle i \rangle\, \restr{t}{U}$.  The boundary conditions
are preserved because restriction commutes with substitution.

(2)--(3) follow from (1) and the definition of composition/inversion
via $\hfilling$, since restriction distributes over $\hfilling$:
$\restr{(\hfilling^i\, A\, [\varphi \mapsto u]\, u_0)}{U} \equiv
\hfilling^i\, \restr{A}{U}\, [\varphi \mapsto \restr{u}{U}]\,
\restr{u_0}{U}$.

(4) is analogous, using the distributivity of restriction over
$\transp$.
\end{proof}

\begin{theorem}[Fiberwise Univalence]
\label{thm:fiberwise-ua}
If $e_\alpha : \restr{A}{U_\alpha} \simeq \restr{B}{U_\alpha}$ is a
family of fiberwise equivalences that is compatible on overlaps
(i.e., $\restr{e_\alpha}{U_{\alpha\beta}} = \restr{e_\beta}{U_{\alpha\beta}}$),
then the fiberwise $\mathsf{ua}$ paths glue to give a global path
$\Path_{\mathcal{U}_i^{\site}}\, A\, B$, provided $\Hone(\covers,
\underline{\mathrm{Aut}}_{\mathcal{U}_i}) = 0$.
\end{theorem}

\begin{proof}
Each $e_\alpha$ gives a path $\mathsf{ua}(e_\alpha) :
\Path\, \restr{A}{U_\alpha}\, \restr{B}{U_\alpha}$ by
Theorem~\ref{thm:univalence}.  The compatibility on overlaps ensures
that the restrictions of these paths agree:
$\restr{\mathsf{ua}(e_\alpha)}{U_{\alpha\beta}} \equiv
\restr{\mathsf{ua}(e_\beta)}{U_{\alpha\beta}}$ (by univalence applied
to the equal equivalences).  This is exactly descent data for the path
type $\Path\, A\, B$ viewed as a sheaf type.  By the descent rule
(Definition~\ref{def:descent-rule}), the local paths glue if and only
if $\Hone = 0$.
\end{proof}

This theorem encapsulates the deepest interaction between cubical and
cohomological reasoning: univalence is \emph{local}, and descent
globalizes it.  No other type theory provides this combination.
